\section{Invariant and arithmetic notation in the family marking}
\label{inv:notation-dictionary}

\paragraph{Base, fibres, and the two peripheral orientations.}
We use $B=\mathbb P^1_t$, $p_1=0$, $p_2=1$, $p_0=\infty$,
$\Sigma=\{p_0,p_1,p_2\}$, and $D=W+S_1+S_2$ with reduced
components. The full fibres are $f^*p_1=3S_1$, $f^*p_2=4S_2$,
and $f^*p_0=W$. A positive base loop is counterclockwise in a
holomorphic local coordinate. The analytic root deck generator is
clockwise, $s_j(g_jz)=e^{-2\pi i/m_j}s_j(z)$, with
$g_1g_2g_0=1$. The lifted positive base loop ends at the inverse deck
translate; its endpoint identification applies $g_j$. Thus its
homology action is $A_j$, not $A_j^{-1}$.
The clockwise boundary orientation used in the collar integration
is stated separately in \cref{inv:C3.2}. References there to a
peripheral generator refer to this fixed analytic generator.

\paragraph{Integral marking and coordinate order.}
The homology basis is
$(\hat\gamma,\hat u,\hat w,\hat\delta)$, and
$(\gamma,u,w,\delta)$ denotes its dual cohomology basis.
The letter $w$ in that cohomology basis is a covector, whereas the
scalar family parameter below is $e^{-\pi ic}$.
In the ordered additive coordinates $(\zeta_1,\zeta_2)$,
\[
 \Pi_c=\begin{pmatrix}6\mu&\tau&1&0\\
             \beta_c&\mu&0&1\end{pmatrix},
 \qquad \beta_c=\beta^{\rm part}+c,
 \qquad (\beta^{\rm part}+\tau)(p_0)=0.
\]
Its rows are $\sigma_1=6\mu\gamma+\tau u+w$ and
$\sigma_2=\beta_c\gamma+\mu u+\delta$.
The actual monodromy matrices in this column basis are
\[
 A_1=\begin{pmatrix}1&0&0&0\\6&0&1&0\\-6&-1&-1&0\\-2&1&0&1\end{pmatrix},
 \quad
 A_2=\begin{pmatrix}1&0&0&0\\0&0&-1&0\\-6&1&0&0\\3&0&1&1\end{pmatrix},
 \quad A_0=(A_1A_2)^{-1}.
\]
No row interchange is implicit: if a local calculation orders the
coordinates as $(\zeta_2,\zeta_1)$, its period matrix is
$J\Pi_c$, $J=\left(\begin{smallmatrix}0&1\\1&0\end{smallmatrix}\right)$,
and its Hermitian coefficient matrix is $JH J^{\mathsf t}$.
Both the main invariant proofs and the retained working-edition
fibre formulas use the displayed $(\zeta_1,\zeta_2)$ order.

\paragraph{The exact Engel and reversed-meridian transitions.}
Engel's base indices $(0,1,\infty)$ correspond to $(1,2,0)$ here.
Writing his common integral basis as column coordinates, the family
basis is given by the columns of
\[
 P_{\rm mark}=\begin{pmatrix}
 0&-2&-1&-2\\2&-1&0&-1\\0&2&1&3\\-1&0&0&0
 \end{pmatrix},\qquad
 T_0P_{\rm mark}=P_{\rm mark}A_1,\quad
 T_1P_{\rm mark}=P_{\rm mark}A_2.
\]
The invariant source's reversed-meridian matrix is
\[
 Q=P_{\rm mark}S,\qquad
 S=\begin{pmatrix}-1&0&0&0\\-6&0&-1&0\\6&-1&0&0\\2&1&1&1\end{pmatrix},
 \qquad
 Q=\begin{pmatrix}2&-1&0&-2\\2&-1&0&-1\\0&2&1&3\\1&0&0&0\end{pmatrix}.
\]
It satisfies $T_0Q=QA_1^{-1}$ and $T_1Q=QA_2^{-1}$.
This is the explicit convention change of \cref{inv:B1.0}.
For invariant finite twist vectors $v_1,v_2$, Engel's positive-root
affine vectors are $a_0=-P_{\rm mark}v_1/3$ and
$a_1=-P_{\rm mark}v_2/4$; his distinguished cycle is
$P_{\rm mark}\hat\delta$.

\paragraph{Multiplicative parameters and special-fibre coordinates.}
Set $v=e^{2\pi i\zeta_2}$. The first period acts by
$(\zeta_1,v)\mapsto(\zeta_1+6\mu,e^{2\pi i\beta_c}v)$.
Engel's contracting generator is its inverse over $-6P$, on the
complement of the zero section of $\mathcal O(P-O)$.
With the reference inverse lift transported from $c=0$,
\[
 w=e^{-\pi ic},\qquad u_{\rm Engel}=w^2,
 \qquad 0<|w|<216,\quad 0<|u_{\rm Engel}|<46656.
\]
A different reference lift multiplies $u_{\rm Engel}$ by a fixed
$\lambda_0\in\mathbb C^*$. Thus $u_{\rm Engel}$ is distinct from
the integral covector $u$ and from an elliptic logarithm denoted $u$
in the period-transcendence section.
The special values use $\rho=e^{\pi i/3}$ rather than
$e^{2\pi i/3}=\rho-1$; the two conventions differ by the modular
translation $\tau\mapsto\tau+1$. The Albanese parameters and their
upper-half-plane representatives are
\[
 \begin{array}{c|c|c|c}
 &\kappa&\text{upper-half-plane period}&\text{nome}\\\hline
 S_1&c+\frac12-\frac{3i}{2\pi}\log48&-\kappa_1&
             q_1=e^{-2\pi i\kappa_1}=-w^2/48^3\\
 S_2&c+\frac32-\frac{3i}{\pi}\log6&-\kappa_2/2&
             q_2=e^{-\pi i\kappa_2}=iw/216.
 \end{array}
\]
The factor $1/2$ is the full order-four Albanese lattice projection.
In particular $q_1=(27/64)q_2^2$.

\paragraph{Determinants, orientations, and cohomology generators.}
Write $y=\Im\tau>0$, $m=\Im\mu$, $r=\Im\beta_c$, and
$D_c=r-6m^2/y=D_{\rm part}+\Im c$. The symbol $\Delta$ used
for this real quantity in the invariant fibre formulas means $D_c$,
not a Weierstrass discriminant. In admissible fibres $D_c<0$.
The real period determinant in the row order
$(\Re\zeta_1,\Im\zeta_1,\Re\zeta_2,\Im\zeta_2)$ is $yD_c$.
Consequently the marked real orientation is opposite to the complex
orientation, and for
$\eta_F=u\wedge w+6\gamma\wedge\delta$,
\[
 \int_{F,\operatorname{or}_{\mathbb C}}\eta_F^2=-12,
 \qquad
 \Gamma=\frac{i}{2D_c}
       \left(d\bar\zeta_2-\frac myd\bar\zeta_1\right)
       \quad\text{represents }\gamma^{0,1}.
\]
The subscript on $\eta_F$ separates this integral fibre class from
Dedekind's eta function and from Weierstrass quasi-periods.
The primitive global class is $\mathfrak g=g_\gamma$.
For standard data $(l_0,l_1,l_2)=(0,1,-1)$ and
$p=12l_0-4l_1-3l_2=-1$,
\[
 \exp(12\mathfrak g)=f^*\mathcal O_B(1),\quad
 \mathcal O_X(S_1-S_2)=\exp(\mathfrak g),\quad
 [G_0]=12\mathfrak g.
\]
In general the first equality is
$\exp(12g_\gamma)=f^*\mathcal O_B(-p)$.
Here $\exp$ means the sheaf map $z\mapsto e^{2\pi iz}$,
$\int_B\omega_{FS}=1$, $c_{1,BC}(L)=[iF_L/(2\pi)]$, and the
unnormalized Chern--Ricci form represents $2\pi c_{1,BC}(T_X)$.
Thus the prescribed Chern--Ricci form is $\pi f^*\omega_{FS}$.

\paragraph{Root and coarse line bundles.}
For the root curve $\mathcal B$ of orders $(3,4)$ put
$R=R_1R_2^{-1}$, so
$R_1=R^4$, $R_2=R^3$, $\rho^*\mathcal O_B(1)=R^{12}$,
and $K_{\mathcal B}=R^{-7}$.
The notation $\mathcal O_{\mathcal B}(n)$ means $R^n$;
$\mathcal O_B(n)$ is the ordinary coarse degree-$n$ line bundle.
They have different degrees and stabilizer characters.
The logarithmic sheaves are
$\mathcal L=\Omega_B^1(\log\Sigma)$,
$\mathcal A=\Omega_X^1(\log D)$, and
$\mathcal E=\mathcal A/f^*\mathcal L$.
Every coherent base-change assertion uses the full scheme fibre.
The finite cyclic covers and subsequent invariant functor retain its
nilpotent normal directions.

\paragraph{Theta norm, modular coordinate, and heights.}
The theta variable has periods $(1,\tau)$; software using a first
argument of period $\pi$ receives $\pi z$. With $q=e^{2\pi i\tau}$,
\[
 N(z,\tau)=\left|\frac{\vartheta_1(z,\tau)}{\eta(\tau)}\right|
       e^{-\pi(\Im z)^2/\Im\tau},\qquad
 D_{\rm part}=\frac6\pi\log\frac{N(\mu,\tau)}{2\sqrt3},
 \qquad t=\frac{j(\tau)}{1728}.
\]
The modular discriminant is
$\Delta_0=q\prod_{n\ge1}(1-q^n)^{24}$, so
$t=E_4^3/(1728\Delta_0)$ and $t-1=E_6^2/(1728\Delta_0)$.
For the rational family
\[
 E_t:Y^2=X^3-\tfrac34t^3(t-1)X+\tfrac14t^4(t-1)^2,
 \qquad P_t=(0,t^2(t-1)/2),
\]
the invariant differential is $dX/(2Y)$ and
$\Delta_E=27t^8(t-1)^3$.
In the lattice model $X=\wp(u)$ and $Y=\wp'(u)/2$;
the lattice discriminant is
$g_2^3-27g_3^2=(2\pi)^{12}\eta(\tau)^{24}$ when the
period lattice is $\mathbb Z+\tau\mathbb Z$.
We use the height of $(O)$,
$\widehat h_O(P)=\frac12\lim n^{-2}h(x(nP))$.
PARI's \texttt{ellheight}, the surface self-pairing, and the
rank-one BSD regulator use $2\widehat h_O$.
The full real period is $\Omega=\int_{E(\mathbb R)}|\omega|$,
including both components if disconnected.
The archimedean local height is exactly $-\log N$.
At finite places $|p|_v=p^{-1}$ and $n_v=[K_v:\mathbb Q_p]$;
complex embeddings are counted separately in the global sum.

For period arguments, $\eta_j$ is defined by
$\zeta(z+\omega_j)=\zeta(z)+\eta_j$, with
$\omega_1\eta_2-\omega_2\eta_1=-2\pi i$.
The real-linear extension of these quasi-periods is denoted $\eta(u)$.
In the critical-point proof, the Cayley coordinate is
$\chi=(\tau-i)/(\tau+i)$, $z=k\chi^2$, and
$\beta_{\rm cusp}=\log(5+2\sqrt6)/(2\pi)$;
neither symbol is the family parameter $w$ or the period $\beta_c$.

\paragraph{Shared statements.}
The maximum has the single label \cref{fam:src:C.2.3}, with aliases
\texttt{ari:C.8.5} and \texttt{push:src:E.4.3}; the independent
interval, duplication, and radial arguments are all retained.
Tangent cohomology, the parameter Kodaira--Spencer map, and the
one-dimensional Kuranishi germ are collected in
\cref{fam:src:A.5.2}, with aliases
\texttt{inv:D1.3}, \texttt{inv:D2.1}, and \texttt{inv:D3.1}.
The exact parameter classification is \cref{fam:src:A.5.1};
the canonical-pencil statement is \cref{fam:src:A.1.1}.
The working Hodge-line proof is \cref{mapsI:hodge:line},
and the finite normal-order calculation is
\cref{wb:ffcert:thm:normal-trivialization}.
Working structure-sheaf and cusp statements share \cref{inv:A1.0};
working Picard restriction shares \cref{inv:C3.1};
the global fibre-class extension shares \cref{inv:A3.2};
the Bott--Chern multiplication statement shares \cref{inv:C3.3}.
